%======================================================================
%  Literature review: the state of the art
%======================================================================

\section{The state of the art}
\label{sec:literature}

Two bodies of work bear on this thesis, and they have grown up almost
independently of each other.  One asks whether a general-purpose CPU can
run a language model quickly enough to be worth using, and answers it by
rewriting arithmetic: quantized weights, cheaper kernels, better use of
the matrix units.  The other asks how much performance depends on
\emph{where} data sits in a tiered memory system, and answers it with
placement experiments, historically on HPC codes and more recently on
the key--value cache of a transformer.  Each line is mature.  \textbf{Their intersection, on a CPU that
actually has two memory tiers under software control, is close to
empty}, and that emptiness is what this thesis is built on.

\subsection{Scope of the review}

The review covers three bodies of work, chosen because each constrains
the study in a different way.  The first is the line of research that
made CPU inference practical through quantization, kernel design and use
of the matrix units.  It supplies the performance figures any result
here will be read against.  The second is work on \emph{memory tiering},
in high-performance computing where it is mature and in transformer
inference where it is not, which supplies the mechanism this thesis
tests.  The third is the smaller methodological literature on evaluating
inference systems, which constrains how the measurements must be taken
and reported.

Two boundaries are worth stating.  Work on training, on distributed
serving and on accelerator microarchitecture is excluded: it is a large
literature, and none of it bears on a single-socket placement question.
And a number of the works discussed below are preprints rather than
refereed publications.  They are cited where they establish what the
field is currently attempting, and their status is noted where a
quantitative claim rests on them, since an unrefereed number is a weaker
foundation than a refereed one.

\secref{sec:related-work} then departs from this survey and reads four
sources closely, on the grounds that they alone were measured on the
processor family studied here.

\subsection{Making the CPU a credible inference target}
\label{sec:lit-arithmetic}

CPU inference became a serious proposition when weight-only quantization
was combined with kernels written for it.  Shen et al.\ paired an
automatic INT4 flow with a purpose-built runtime and reported 20--80\,ms
per generated token for 6--20\,B models on a single fourth-generation
Xeon socket, within a percent of FP32 accuracy~\cite{shen2023efficientcpu}.
The significance is less the latency than the demonstration that the
gap to an accelerator was a software gap as much as a hardware one.

In practice most CPU inference now happens through
\lcpp{}~\cite{llamacpp2026}, the runtime adopted here, and the published
measurements of it provide a useful calibration.  Malakhov benchmarks seven 4-bit GGUF
models across three CPU environments, including a Xeon Platinum 8592+,
reporting roughly 120 tokens/s for a 1\,B model, 45 for a 7\,B model and
15 for a 24\,B model~\cite{malakhov2025medlocalgpt}.  That machine is
the DDR-only sibling of the processor studied here, being the same
generation with a comparable core count and no on-package HBM.  That
makes it the closest thing in the literature to a published control for the comparison in
\secref{sec:research-gap}.  Kurt supplies the orthogonal sweep,
evaluating 3- to 8-bit K-quant and legacy GGUF formats on one model for
accuracy, perplexity and CPU prefill and decode throughput together,
which is the evidence needed to choose a quantization level on grounds
other than habit~\cite{kurt2026quantization}.

Alongside the runtimes sits a line of work on the kernels themselves,
unified by the observation that a quantized matrix--vector product
spends most of its time on operations the hardware is bad at.  Zhang et
al.\ replace multiply--accumulate in attention with in-register SIMD
lookups and roughly double the speed of a 4-bit LLaMA-7B at 16\,k
context~\cite{zhang2024nomad}.  Wei et al.\ generalise the idea to the
weight path, computing mixed-precision GEMM by table lookup with no
dequantization step at all, and report up to 4$\times$ over \lcpp{}
together with a 70\% reduction in energy~\cite{wei2025tmac}.  Zuo et
al.\ push furthest, fusing eight sub-GEMVs of a ternary layer into a
single AVX-512 loop of masked additions and subtractions, with no
multiplications anywhere in the inner
loop~\cite{zuo2026fairyfuse}.

\begin{table}[htbp]
\centering
\small
\tabsetup
\caption{Arithmetic-side work on CPU inference.  Reported gains are
against the baseline named in each paper and on that paper's hardware;
they are not mutually comparable, and are listed to show the shape of
the field rather than to rank it.}
\label{tab:lit-arithmetic}
\begin{tabularx}{\textwidth}{@{}P{2.5cm}P{3.1cm}YP{3.1cm}@{}}
\toprule
\headrow
\textbf{Work} & \textbf{Platform} & \textbf{Lever} &
\textbf{Baseline and gain} \\
\midrule
Shen et al.\ 2023~\cite{shen2023efficientcpu}
& 4th-gen Xeon, one socket
& INT4 weight-only quantization with a matched runtime
& 20--80\,ms/token for 6--20\,B; \(\leq\)1\% accuracy loss vs FP32 \\
\addlinespace[2pt]
Zhang et al.\ 2024~\cite{zhang2024nomad}
& x86 with SIMD registers
& Attention by in-register lookup instead of multiply--add
& \(\approx\)2$\times$ on 4-bit LLaMA-7B at 16\,k context \\
\addlinespace[2pt]
Wei et al.\ 2025~\cite{wei2025tmac}
& Apple M2-Ultra, Raspberry Pi 5
& Table-lookup mixed-precision GEMM, no dequantization
& Up to 4$\times$ over \lcpp{}; 70\% less energy \\
\addlinespace[2pt]
Zuo et al.\ 2026~\cite{zuo2026fairyfuse}
& Intel Xeon 8558P
& Ternary weights, eight sub-GEMVs fused into one AVX-512 loop
& 29.6$\times$ at kernel level; \textbf{1.24$\times$ end to end} over
  \lcpp{} Q4\_K\_M \\
\bottomrule
\end{tabularx}
\end{table}

\subsection{Where the arithmetic-side lever runs out}
\label{sec:lit-pivot}

The last row of \tabref{tab:lit-arithmetic} is the most informative
number in this section, and it is worth dwelling on.  Zuo et al.\ make
their kernel almost thirty times faster and the model runs 1.24 times
faster.  Nothing went wrong in their engineering: \emph{the arithmetic simply
stopped being what the program was waiting for}.  Once a phase spends
most of its time waiting for weights to arrive, removing instructions
from that phase buys very little, and the marginal return on further
kernel work collapses.

This is the pivot on which the rest of the review turns.  A decade of
effort has moved CPU inference from implausible to routine by making the
work cheaper, and that effort has been successful enough to expose what
lies underneath it.  The interesting question is no longer how few
operations a token can be generated with, but how quickly the bytes
those operations consume can be delivered, which is a question about the
memory system and therefore about placement.
\figref{fig:two-levers} sets the two lines of attack against each other
and marks the point at which they stop being interchangeable.

\begin{figure}[htbp]
\centering
\begin{tikzpicture}[
  font=\small,
  lever/.style={draw=black!65, fill=neutralcol, rounded corners=3pt,
                minimum width=6.1cm, minimum height=1.0cm, align=center},
  work/.style={font=\footnotesize, text width=5.9cm, align=center,
               inner sep=2pt},
  target/.style={draw=black!70, fill=white, rounded corners=3pt,
                 minimum width=6.6cm, minimum height=1.15cm,
                 align=center},
  arr/.style={-{Stealth}, thick, black!60}
]
  \node[lever, fill=ddrcol!12, draw=ddrcol!80!black] (a) at (-3.6,2.6)
    {\textbf{Reduce the work}\\[1pt]
     \footnotesize fewer or cheaper operations per token};
  \node[work, below=0.15cm of a, text=black!70]
    {quantization~\cite{shen2023efficientcpu,kurt2026quantization},
     lookup kernels~\cite{zhang2024nomad,wei2025tmac},
     ternary fusion~\cite{zuo2026fairyfuse}};

  \node[lever, fill=hbmcol!14, draw=hbmcol!80!black] (b) at (3.6,2.6)
    {\textbf{Reduce the wait}\\[1pt]
     \footnotesize fewer or faster bytes per token};
  \node[work, below=0.15cm of b, text=black!70]
    {tier placement~\cite{fang2025kvplacement,jang2026itme},
     cache residency~\cite{zhang2026cacheresident},
     bandwidth-aware mapping~\cite{lu2026thinfer}};

  \node[target] (t) at (0,-0.75)
    {\(T_{\mathrm{step}} \gtrsim
       \max\!\left(\dfrac{F_{\mathrm{step}}}{P_{\mathrm{eff}}},\;
                   \dfrac{D_{\mathrm{step}}}{B_{\mathrm{eff}}}\right)\)};

  \draw[arr, ddrcol!80!black] (-3.6,1.35) -- (-1.4,-0.35);
  \draw[arr, hbmcol!80!black] (3.6,1.35) -- (1.4,-0.35);

  \node[font=\footnotesize, text width=11.4cm, align=center,
        below=0.35cm of t, text=black!75]
    {When the second term dominates, the left lever saturates: Zuo et
     al.\ report 29.6$\times$ at the kernel and 1.24$\times$ end to
     end~\cite{zuo2026fairyfuse}.  This thesis pulls the right lever on
     hardware where both tiers are directly addressable.};
\end{tikzpicture}
\caption{The two ways of shortening a decode step, and the point at
which they stop being interchangeable.  Author's synthesis; the
lower-bound expression is the step model developed in
\secref{sec:roofline}.}
\label{fig:two-levers}
\end{figure}


\subsection{Moving the data instead}
\label{sec:lit-placement}

\subsubsection{What HPC already knows}

Tiered main memory is not new, and the HPC community learned most of the
relevant lessons on Knights Landing.  Peng et al.\ compared MCDRAM with
DRAM across scientific and analytics workloads and found a split that
has aged well: codes with regular access patterns gained up to
3$\times$, while codes with irregular patterns were latency-bound and
could run \emph{slower} in the high-bandwidth tier
alone~\cite{peng2017hybridmemory}.  Their remedy, more hardware threads to keep enough requests in flight
for the bandwidth to be reachable, is the same concurrency argument
developed in \secref{sec:littles-law}, and it is the most direct published
warning that a fast tier is not automatically a faster program.

On the processor used here, Reguly provides the closest non-LLM
reference point, reporting 2.0--4.3$\times$ over the previous generation
on bandwidth-sensitive HPC proxies and, importantly, observing that for
some codes the bottleneck moves from bandwidth to communication
latency once bandwidth stops being scarce~\cite{reguly2023xeonmax}.  That
observation is the reason this thesis does not expect an inference
speedup equal to the STREAM ratio, and it sets the scale against which a
smaller measured effect should be judged.

\subsubsection{Tiering for inference}

Applying the same idea to transformers is recent and, so far, mostly
indirect.  Fang et al.\ give the problem its formal treatment, casting
KV-cache placement across a fast and a slow tier as latency
minimisation and deriving an upper bound of up to 5.87$\times$ over
static placement~\cite{fang2025kvplacement}.  Three properties of that
work matter here.  It is evaluated in a behavioural simulator rather
than on hardware; it covers the decode stage only; and it fixes model
weights in the fast tier by assumption, on the grounds that weights are
read at every step, leaving KV entries as the only variable.  The last
of these is the assumption a Xeon Max socket is able to break, because
64\,GB of HBM, and only 16\,GB per domain under SNC4, is not comfortably
larger than a quantized model of interest.

Neighbouring systems work approaches the same architecture from other
directions.  Jang et al.\ exploit the predictable access pattern of
weights and prefix state to expand KV capacity across host DRAM,
CXL-attached memory and NVMe, reporting a 35.7\% throughput
improvement~\cite{jang2026itme}, and a companion design places a
three-tier stack behind a CXL interface for multi-turn
serving~\cite{jang2026hymcache}.  Most directly comparable to this
thesis, Zhang et al.\ keep reusable weights resident in a GB-scale
last-level cache on a multi-socket CPU, separating weight-centric
operators from attention and KV management, and report 2.04--11.51$\times$
on time per output token against \lcpp{} with equivalent
resources~\cite{zhang2026cacheresident}.  Their argument is this thesis's argument with a different fast tier:
\emph{SRAM instead of HBM, bought with a redesigned runtime rather than
with placement}.

\tabref{tab:lit-tiering} collects this work and makes the pattern
visible.  At larger scale, Lu et al.\ build an inference framework for
the MT-3000
many-core processor of the Tianhe supercomputer, framing the problem
explicitly as maximising locality under a bandwidth constraint and
beating DeepSpeed on 2$\times$V100S by 62--73\% for a 7\,B
model~\cite{lu2026thinfer}.  Chen et al.\ show that the CPU side of a
hybrid deployment is itself the bottleneck, using AMX-specialised
kernels and asynchronous scheduling to obtain 4.62--19.74$\times$
prefill and 1.25--4.09$\times$ decode speedups on large
MoE models~\cite{chen2025ktransformers}; He et al.\ establish that
scaling CPU inference beyond one socket is ordinary
practice~\cite{he2024distributedcpu}.  Taken together these say that
CPU-side memory behaviour is worth improving even when a GPU is
present.

\begin{table}[htbp]
\centering
\small
\tabsetup
\caption{Placement and tiering work relevant to the questions in \secref{sec:research-gap}.  The
final column is the one that matters: almost all inference-side evidence
is simulated, accelerator-attached, or obtained by rebuilding the
runtime.}
\label{tab:lit-tiering}
\begin{tabularx}{\textwidth}{@{}P{2.45cm}P{3.3cm}YP{2.7cm}@{}}
\toprule
\headrow
\textbf{Work} & \textbf{Tiers} & \textbf{What is placed} &
\textbf{Evidence} \\
\midrule
Peng et al.\ 2017~\cite{peng2017hybridmemory}
& MCDRAM / DRAM, one socket
& Whole application; regular vs irregular access compared
& Measured, HPC codes \\
\addlinespace[2pt]
Reguly 2023~\cite{reguly2023xeonmax}
& HBM2e / DDR5, Xeon Max
& Whole application
& Measured, HPC proxies \\
\addlinespace[2pt]
Fang et al.\ 2025~\cite{fang2025kvplacement}
& GPU HBM3 / CPU LPDDR5X over NVLink-C2C
& KV entries only; weights pinned in HBM by assumption
& \textbf{Simulated}, decode only \\
\addlinespace[2pt]
Jang et al.\ 2026~\cite{jang2026itme}
& Host DRAM / CXL / NVMe
& Weights and prefix cache, by access predictability
& Measured, capacity expansion \\
\addlinespace[2pt]
Zhang et al.\ 2026~\cite{zhang2026cacheresident}
& LLC / DRAM, multi-socket CPU
& Weights held cache-resident; KV scaled separately
& Measured, \textbf{requires a redesigned runtime} \\
\addlinespace[2pt]
\textbf{This thesis}
& \textbf{HBM2e / DDR5, one Xeon Max socket}
& \textbf{Whole process first; then weights against KV state}
& \textbf{Measured, unmodified runtime, residency verified} \\
\bottomrule
\end{tabularx}
\end{table}


\subsection{Whether decode is bandwidth-bound, and why it might not be}
\label{sec:lit-bandwidth}

The claim that autoregressive decode is memory-bound is close to folklore
now, and it is usually justified with a roofline argument.  Yuan et al.\
supply the canonical version, building a survey of inference efficiency
on an explicit roofline analysis and showing why low-batch decode sits
far to the left of the ridge point~\cite{yuan2024roofline}.  Bi et al.\
turn the same model into a benchmarking framework and add a detail this
thesis needs: operational intensity is not a fixed property of a model
but moves with sequence length and depth~\cite{bi2026rooflinebench},
which is precisely the mechanism behind the capacity sweep in
\secref{sec:working-set}.

What the folklore leaves out is that being memory-bound and being
bandwidth-limited are different conditions, and only the second is cured
by a faster tier.  Chen reports a cross-GPU study of batch-1 decode in
which faster memory conspicuously fails to deliver proportional latency
gains: an L4 reaches 81\% of its peak bandwidth while an H100 reaches
27\%, and launch overhead and kernel choice matter more than the
bandwidth on the specification sheet~\cite{chen2026memorybound}.  The
work is GPU-side, so it transfers as a mechanism rather than as a
comparable measurement, but the mechanism is exactly the one that would
produce a disappointing HBM result here, and it is the strongest
published reason to state \(H_{1b}\), that the observed ratio will fall
short of the STREAM ratio, as a hypothesis rather than as a concession.

Arif et al.\ add a third regime.  Studying models from 8\,B to 671\,B on
GPU clusters, they describe reasoning-heavy workloads as
\emph{capacity-bound} rather than compute- or bandwidth-bound, and
identify KV-cache fragmentation as a mechanism that throttles
throughput before any bandwidth ceiling is
reached~\cite{arif2026inferencescaling}.  For this thesis that is both a
vocabulary and a warning: the transition
\secref{sec:research-gap} sets out to locate could be produced by
fragmentation rather than by tier capacity, and the two have to be told
apart.

\subsection{How the field measures, and how it gets it wrong}
\label{sec:lit-method}

A memory-placement result is a comparison, and comparisons of inference
systems are easy to draw badly.  Agrawal et al.\ catalogue the failure
modes across three axes: baselines that conflate engineering effort with
algorithmic novelty, workloads that do not resemble anything deployed,
and normalised metrics that conceal generation stalls.  They also supply
a checklist for avoiding them~\cite{agrawal2025evaluating}.
Their argument is the reason the comparison at the centre of this thesis
holds the build, thread count, model, quantization and prompt fixed and
varies only the memory tier, and the reason inter-token latency will be
reported as a distribution rather than as a mean.

Two further papers shape the protocol.  Salaria et al.\ propose metrics
for the cost of benchmarking itself and show that a well-chosen subset
of a sweep can retain most of its information.  Reducing generation
length from 1024 to 128 tokens preserved 96.6\% of the accuracy of the
full measurement at 2.7$\times$ less cost~\cite{salaria2025metametrics}.
On a shared allocation with a 24-hour limit that is not a convenience
but a design constraint, and it gives a defensible basis for a reduced
factor matrix instead of an arbitrary one.  Stuhlmann et al.\ report
latency, throughput, energy and startup time side by side across
engines and quantizations and conclude that no configuration wins
everywhere~\cite{stuhlmann2025bench360}; their results are GPU-side, but
the reporting structure is the right template for a results chapter
whose honest conclusion will also be conditional.

For placing CPU inference in a wider hardware context without expanding
the scope of the experiments, Li et al.\ provide a cross-platform survey
covering CPU, GPU, FPGA, ASIC and processing-in-memory designs, with
tokens per second and tokens per joule at small batch
sizes~\cite{li2024hardwareperspective}.

\subsection{Where this thesis sits}
\label{sec:lit-position}

Reading the two lines of work together produces a fairly sharp picture.
The arithmetic-side line has largely done its job and is now limited by
data movement.  The memory-side line has strong measured results in HPC,
where the object being placed is an array, and much weaker measured
results in inference, where the object being placed is a weight tensor
or a KV cache.  The inference-side tiering evidence is dominated by
simulation~\cite{fang2025kvplacement}, by accelerator-attached
hierarchies reached over an interconnect~\cite{jang2026itme,jang2026hymcache},
or by results that required rebuilding the runtime around a different
fast tier~\cite{zhang2026cacheresident}.

A Xeon Max socket is unusual in that it removes all three of those
qualifications at once.  Both tiers are directly load/store addressable
by the same cores, so placement is a page policy rather than a transfer;
the runtime does not have to be modified for the primary comparison; and
the fast tier is small enough relative to a quantized model that the
weights-fit assumption underlying the published formulation can be made
to fail on purpose.  What is missing from the literature, and what
\secref{sec:research-gap} sets out to supply, is the corresponding
measurement: a phase-resolved comparison of verified local-HBM and
local-DDR placement for quantized \lcpp{} inference, with page residency
and achieved bandwidth reported, on hardware rather than in a model.

Whether that measurement is the first of its kind is not a claim this
review makes.  The search described above was not systematic enough to
support it, and the more useful framing is in any case comparative: the
prediction exists, the bound has been derived, and what has not been
done is to hold it against a real machine.
